\documentclass{article}
\usepackage{amsmath, amssymb, amsthm}
\usepackage{hyperref}
\usepackage{geometry}
\geometry{margin=1in}

\title{Erd\H{o}s Problem \#836}
\date{Last updated: 2025-08-31}
\author{Source: erdosproblems.com}

\begin{document}
\maketitle

\noindent\textbf{Status:} OPEN \quad \textbf{No prize}

\noindent\textbf{Tags:} graph theory, hypergraphs, chromatic number

\medskip

\noindent\textbf{Problem Statement:}

Let $r\geq 2$ and $G$ be a $r$-uniform hypergraph with chromatic number $3$ (that is, there is a $3$-colouring of the vertices of $G$ such that no edge is monochromatic). Suppose any two edges of $G$ have a non-empty intersection. Must $G$ contain $O(r^2)$ many vertices? Must there be two edges which meet in $\gg r$ many vertices?




A problem of Erd\H{o}s and Shelah. The Fano geometry gives an example where there are no two edges which meet in $r-1$ vertices. Are there any other examples? Erd\H{o}s and Lov\'{a}sz \cite{ErLo75} proved that there must be two edges which meet in $\gg \frac{r}{\log r}$ many vertices. Alon has provided the following counterexample to the first question: as vertices take two sets $X$ and $Y$ of sizes $2r-2$ and $\frac{1}{2}\binom{2r-2}{r-1}$ respectively, where $Y$ corresponds to all partitions of $X$ into two equal parts. The edges are all subsets of $X$ of size $r$, and also all sets consisting of a subset of $X$ of size $r-1$ together with the unique element of $Y$ corresponding to the induced partition of $X$. This hypergraph is intersecting, its chromatic number is $3$, and it has $\asymp 4^r/\sqrt{r}$ many vertices.




References

[ErLo75] Erd\H{o}s, P. and Lov\'{a}sz, L., Problems and results on {$3$}-chromatic hypergraphs and some related questions . (1975), 609--627.


\bigskip
\noindent\small{Source: \url{https://www.erdosproblems.com/836}}
\end{document}
